\chapter{Forward Dynamics --- Inertia Matrix Methods}
\label{chap:6}
Forward dynamics is the problem of finding the acceleration of a rigid-body system in response to given applied forces. It is used mainly for simulation; and it is sometimes called ‘direct dynamics’, or simply ‘dynamics’. In this chapter and the next, we examine the forward dynamics of kinematic trees. The dynamics of closed-loop systems is covered in Chapter 8. There are two main approaches to solving the forward dynamics problem for a kinematic tree: 1. form an equation of motion for the whole system, and solve it for the acceleration variables; or 2. propagate constraints from one body to the next in such a way that the accelerations can be calculated one joint at a time. We examine the first approach in this chapter, and the second in Chapter 7. The first approach involves calculating the elements of an n × n joint-space inertia matrix, and then factorizing this matrix in order to solve a set of n linear equations for the acceleration variables. Such algorithms have O(n3) complexity in the worst case, and are sometimes referred to collectively as O(n3) algorithms. However, it would be more accurate to describe them as O(nd2) algorithms, where d is the depth of the tree. The second approach involves making a fixed number of passes through the tree, each pass performing a fixed number of calculations per body. Algorithms that take this approach have O(n) complexity, and are therefore the best choice for systems containing a large number of bodies. However, the O(nd2) algorithms can match or slightly exceed the speed of the O(n) algorithms on trees with only a few bodies, or that are branched enough to have a small depth. The O(nd2) algorithms are also an important component in closed-loop dynamics algorithms. The efficiency trade-offbetween these two approaches is discussed in Chapter 10.

This chapter commences with some basic material on how to solve the forward dynamics problem using the joint-space inertia matrix. Then Section 6.2 presents the composite-rigid-body algorithm, which is the fastest algorithm for calculating this matrix; and Section 6.3 presents a second derivation that reflects how the algorithm was originally discovered and described. The next two sections describe the phenomenon of branch-induced sparsity, and present two factorization algorithms to exploit it. This sparsity is the reason why the cost of calculating the joint-space inertia matrix is only O(nd), and the cost of factorizing it only O(nd2). Finally, Section 6.6 presents some background material on the forward dynamics problem. 6.1 The Joint-Space Inertia Matrix The equation of motion for a kinematic tree can be expressed in matrix form as τ = H(q) ¨q + C(q, ˙q, f x) , (6.1) where q, ˙q, ¨q and τ are vectors of joint position, velocity, acceleration and force variables, f x is a vector of external forces (if any), H is the joint-space inertia matrix, and C is the joint-space bias force. H and C are the coefficients of the equation of motion, while τ and ¨q are the variables. This equation shows H as a function of q, and C as a function of q, ˙q and f x. Once these functional dependencies are understood, they are usually omitted. The most useful piece of information they convey is the fact that H depends only on the position variables. If the system is at rest, and there are no forces acting on it other than τ, then C = 0 and the equation of motion simplifies to τ = H ¨q. In the forward dynamics problem, τ is given, and the task is to calculate ¨q. The methods described in this chapter solve the problem by the following procedure: 1. calculate C, 2. calculate H, 3. solve H¨q = τ −C for ¨q. These three steps have computational complexities of O(n), O(nd) and O(nd2), respectively, where d is the depth of the tree. Given that d is bounded by d ≤n, the worst-case complexities for steps 2 and 3 are O(n2) and O(n3). At the other extreme, if d is subject to a fixed upper limit (which can happen in practice) then the complexity of all three steps is O(n). As explained below, step 1 can be accomplished easily by an inverse dynamics algorithm, so this chapter is concerned mostly with steps 2 and 3. H is a symmetric, positive-definite, n × n matrix that is organised into an NB × NB array of submatrices. (Recall that NB = NJ for a tree.) If ni is the degree of freedom of joint i, then Hij is the ni × nj submatrix of H occupying block-row i and block-column j. If pi = Pi k=1 nk then block-row i consists of

rows pi−1 + 1 to pi inclusive, and block-column j consists of columns pj−1 + 1 to pj inclusive. Physically, Hij relates the acceleration at joint j to the force at joint i. If every joint in the tree has only one degree of freedom, then n = NB, and Hij is the 1 × 1 matrix on row i and column j of H. Suppose we have an inverse dynamics calculation function, ID, like the one described in Chapter 5. This function performs the calculation τ = ID(model, q, ˙q, ¨q, f x) . On comparing this expression with Eq. 6.1, it follows immediately that C = ID(model, q, ˙q, 0, f x) (6.2) and H¨q = ID(model, q, ˙q, ¨q, f x) −C . (6.3) Thus, the value of C can be obtained with a single call to ID, and the product of H with any desired vector is the difference between two calls to ID. Let us define a differential inverse dynamics function, IDδ, which calculates the change in joint force corresponding to a change in desired acceleration. This function is given by IDδ(model, q, ¨q) = ID(model, q, ˙q, ¨q, f x) −ID(model, q, ˙q, 0, f x) . (6.4) With this function, we can calculate column α of H as follows: Hδα = IDδ(model, ˙q, δα) , where δα is an n-dimensional coordinate vector having a 1 in its αth element and zeros elsewhere. By repeating this calculation with α taking every value from 1 to n, we can calculate the whole of H one column at a time. This idea is the basis of methods 1 and 2 described in Walker and Orin (1982). There is a great deal of cancellation of terms on the right-hand side of Eq. 6.4. In particular, every term depending on either ˙q or f x cancels out, which is why these vectors have been omitted from the argument list of IDδ. It is therefore possible to formulate a greatly simplified version of the recursive Newton-Euler algorithm specifically to calculate IDδ. The equations and pseudocode for this algorithm are shown in Table 6.1. The pseudocode omits the calculation of Si and iXλ(i) because they are assumed to be available as a by-product of calculating C. The algorithm in Table 6.1 has the advantage of being simple and straightforward. However, it is not the most efficient way to calculate H. The fastest algorithm for that job is the composite-rigid-body algorithm, which is the subject of the next section.

Equations: ai = iXλ(i) aλ(i) + Si ¨qi (a0 = 0) fi = Ii ai + X j∈µ(i) iX∗ j fj τi = ST i fi Algorithm: a0 = 0 for i = 1 to NB do ai = iXλ(i) aλ(i) + Si ¨qi fi = Ii ai end for i = NB to 1 do τi = ST i fi if λ(i)̸ = 0 then fλ(i) = fλ(i) + λ(i)X∗ i fi end end Table 6.1: Equations and algorithm to calculate τ = H ¨q 6.2 The Composite-Rigid-Body Algorithm The kinetic energy of a kinematic tree is the sum of the kinetic energies of its bodies. Thus, T = 1 2 NB X k=1 vT k Ik vk . (6.5) The velocity of body k is given by vk = X i∈κ(k) Si ˙qi , (6.6) where κ(k) is the set of joints that support body k; that is, the set of joints on the path between body k and the base. (See §4.1.4 for descriptions of κ(i), ν(i) and µ(i).) Substituting Eq. 6.6 into 6.5 gives T = 1 2 NB X k=1 X i∈κ(k) X j∈κ(k) ˙qT i ST i Ik Sj ˙qj . (6.7) Now, this equation is a sum over all i, j, k triples in which both joint i and joint j support body k. It can therefore be rearranged into T = 1 2 NB X i=1 NB X j=1 X k∈ν(i)∩ν(j) ˙qT i ST i Ik Sj ˙qj , (6.8) in which ν(i) ∩ν(j) is the intersection of the set of bodies supported by joint i and the set of bodies supported by joint j. Equation 6.8 is therefore a sum

over all i, j, k triples in which body k is supported both by joint i and by joint j, which is the same sum as in Eq. 6.7. The value of ν(i) ∩ν(j) is given by ν(i) ∩ν(j) =

ν(i) if i ∈ν(j) ν(j) if j ∈ν(i) ∅ otherwise (6.9) where ∅denotes the empty set. Now, one of the properties of the joint-space inertia matrix is that the kinetic energy of the tree is given by the expression T = 1 2 ˙qTH ˙q = 1 2 NB X i=1 NB X j=1 ˙qT i Hij ˙qj . (6.10) On comparing this equation with Eq. 6.8, it can be seen that Hij = X k∈ν(i)∩ν(j) ST i Ik Sj . (6.11) Let us now introduce a new quantity, Ic i , which is the inertia of the subtree rooted at body i, treated as a single composite rigid body. (This is where the algorithm gets its name.) Ic i is given by Ic i = X j∈ν(i) Ij ; (6.12) and it can be computed recursively using the formula Ic i = Ii + X j∈µ(i) Ic j . (6.13) Combining Eqs. 6.11, 6.12 and 6.9 produces the following formula for the jointspace inertia matrix: Hij =

ST i Ic i Sj if i ∈ν(j) ST i Ic j Sj if j ∈ν(i) 0 otherwise. (6.14) Equations 6.13 and 6.14 are the basic equations of the composite-rigid-body algorithm for a general kinematic tree. For the special case of an unbranched tree, these equations simplify to Ic i = Ii + Ic i+1 (6.15) and Hij = ST i Ic max(i,j) Sj . (6.16)

Equations in Body Coordinates Equation 6.13 can be expressed in body coordinates as Ic i = Ii + X j∈µ(i) iX∗ j Ic j jXi . (6.17) In principle, Eq. 6.14 can be expressed in body coordinates as Hij =

ST i Ic i iXj Sj if i ∈ν(j) ST i iX∗ j Ic j Sj if j ∈ν(i) 0 otherwise. (6.18) However, to devise an efficient calculation scheme for this equation, it is necessary to define a new quantity, jFi, which is the value of Ic i Si expressed in body j coordinates; that is, jFi = jX∗ i Ic i Si. To calculate H, we need one instance of jFi for every i, j pair satisfying j ∈κ(i). They can be calculated recursively using the formula λ(j)Fi = λ(j)X∗ j jFi , (iFi = Ic i Si) (6.19) with i taking every value in the range 1 to NB, while j progresses (separately for each i) through the values i, λ(i), λ(λ(i)), and so on, back to the base. The joint-space inertia matrix is then given by Hij =

jF T i Sj if i ∈ν(j) HT ji if j ∈ν(i) 0 otherwise. (6.20) Equations 6.17, 6.19 and 6.20 are the equations of the composite rigid body algorithm in body coordinates. Algorithm Details Table 6.2 shows the pseudocode for the body-coordinates version of the algorithm, together with both the basic equations and the equations for the bodycoordinates algorithm. Bear in mind that the symbols Ii, Ic i and Si in the basic version represent either abstract tensors or tensors expressed in an unspecified common coordinate system, whereas the same symbols in the body-coordinates equations represent quantities that are expressed in the coordinate system of body i. The pseudocode does not include code to calculate Si or iXλ(i) because these quantities are assumed to be available as a by-product of applying the recursive Newton-Euler algorithm to calculate C in Eq. 6.1. In implementing Eq. 6.17, we have transformed it from a calculation using µ(i) to a calculation using λ(i). A literal translation of Eq. 6.17 into pseudocode looks like this:

Basic Equations: Ic i = Ii + X j∈µ(i) Ic j Hij =

ST i Ic i Sj if i ∈ν(j) ST i Ic j Sj if j ∈ν(i) 0 otherwise Equations for Body-Coordinates Algorithm: Ic i = Ii + X j∈µ(i) iX∗ j Ic j jXi λ(j)Fi = λ(j)X∗ j jFi (iFi = Ic i Si) Hij =

jF T i Sj if i ∈ν(j) HT ji if j ∈ν(i) 0 otherwise Algorithm: H = 0 for i = 1 to NB do Ic i = Ii end for i = NB to 1 do if λ(i)̸ = 0 then Ic λ(i) = Ic λ(i) + λ(i)X∗ i Ic i iXλ(i) end F = Ic i Si Hii = ST i F j = i while λ(j)̸ = 0 do F = λ(j)X∗ j F j = λ(j) Hij = F TSj Hji = HT ij end end Table 6.2: The composite rigid body equations and algorithm for i = NB to 1 do Ic i = Ii for each j in µ(i) do Ic i = Ic i + iX∗ j Ic j jXi end end However, the same calculation can be performed in a different order, but with the same end result, using the following code fragment (cf. the implementation of Eq. 5.20 in §5.3): for i = 1 to NB do Ic i = Ii end for i = NB to 1 do if λ(i)̸ = 0 then Ic λ(i) = Ic λ(i) + λ(i)X∗ i Ic i iXλ(i) end end The algorithm starts with two preparatory steps: setting H to zero and setting each Ic i to Ii. The first step is included to remind the reader that

body acceleration joint force sα 0 Ii c sα sα joint i Ij c sα ( j ν(i)) 0 Figure 6.1: Kinematic tree undergoing an acceleration of ¨q = δα the rest of the code will initialize only the nonzero submatrices in H. This step may or may not be necessary, depending on how H is to be used. The main loop then implements the calculations for Eqs. 6.17, 6.19 and 6.20. F is a local variable that is initially set to Ic i Si (and therefore equal to iFi in the equations). The while loop transforms this variable successively through the coordinate systems of every body j on the path from body i to the base (i.e., every j ∈κ(i)), as it calculates the submatrices Hij and Hji. Note that setting Hji will only be necessary if the upper triangle is to be accessed. If d is the depth of the kinematic tree, then the while loop will never execute more than d −1 times for any value of i. It therefore follows that the computational cost of the composite-rigid-body algorithm cannot be more than O(nd), and that the number of nonzero elements in H is at most O(nd). A detailed cost analysis appears in Section 10.3.2. 6.3 A Physical Interpretation This section presents an alternative derivation of the composite-rigid-body algorithm that is closer to the original version, the purpose being to show some of the physical insights that played a role in the discovery of this algorithm. As mentioned earlier, if the kinematic tree is at rest, and there are no forces acting on it other than τ, then the equation of motion simplifies to τ = H ¨q. If we set ¨q = δα under these conditions, then τ will be the αth column of H. Figure 6.1 shows what is happening in the tree when ¨q = δα. The αth variable in ¨q must belong to a joint. We therefore identify the joint that owns this variable, and call it joint i. Likewise, we identify the column in Si that corresponds to the αth variable, and call it sα. (sα is a spatial motion vector.) The motion of the tree can now be stated as follows: the acceleration across joint

i is sα, and the acceleration across every other joint is zero. As a consequence, the acceleration of every body in the subtree supported by joint i is sα, and every other body has an acceleration of zero. In other words, aj =  sα if j ∈ν(i) 0 otherwise. (6.21) Let f B j and fj be the net force on body j and the total force transmitted across joint j, respectively. As there are no velocity terms, f B j is given by f B j = Ij aj =  Ij sα if j ∈ν(i) 0 otherwise. (6.22) Now, joint j is the only connection between the subtree ν(j) and the rest of the system. As there are no forces acting on the bodies other than those transmitted across the joints, it follows that fj must be the net force acting on subtree ν(j). However, the net force on a subtree is also the sum of the net forces on the bodies in the subtree, so we have fj = X k∈ν(j) f B k . (6.23) Combining these two equations gives fj =

P k∈ν(j) Ik sα if j ∈ν(i) P k∈ν(i) Ik sα if j ∈κ(i) 0 otherwise, which simplifies to fj =

Ic j sα if j ∈ν(i) Ic i sα if j ∈κ(i) 0 otherwise (6.24) in view of Eq. 6.12. Finally, the joint force variables are given by τj = ST j fj =

ST j Ic j sα if j ∈ν(i) ST j Ic i sα if j ∈κ(i) 0 otherwise. (6.25) These are the subvectors of column α of H. The second case in Eq. 6.24 reveals the two key facts that led to the discovery of the composite-rigid-body algorithm: that fi is the product of a joint motion vector with a composite-rigid-body inertia, and that fj = fi for all j ∈κ(i) (the dark-grey joints in Figure 6.1). The algorithm therefore consists of: 1. calculating the composite-rigid-body inertias via Eq. 6.17; 2. calculating fi = Ic i sα in body i coordinates;

= nonzero submatrix (c) (a) (b) (d) (e) 1 2 3 4 1 4 3 2 Figure 6.2: Examples of branch-induced sparsity 3. transforming fi to the coordinate system of each body in κ(i); and 4. calculating the elements of H in the portion of column α above the main diagonal, and copying them to the portion of row α below the main diagonal, the two being the same because of symmetry. The only difference between this algorithm and the one in Section 6.2 is that the latter processes the columns in blocks, rather than individually, one block being all of the columns pertaining to a single joint. If joint i has ni degrees of freedom, then the quantity iFi in Eq. 6.19 is simply a 6 × ni matrix whose columns are the forces fi = Ic i sα for each α belonging to joint i. 6.4 Branch-Induced Sparsity Equation 6.14 implies that some elements of H will automatically be zero simply because of the connectivity. According to this equation, Hij will be zero whenever i is neither an ancestor of j, nor a descendant, nor equal to j. This can happen only if i and j are on different branches of the tree, so we call it branch-induced sparsity. Figure 6.2 shows some simple connectivity graphs and the sparsity patterns they produce. Example (a) shows an extreme case in which every body is connected directly to the base, resulting in the greatest possible amount of branch-induced sparsity. In this case, every off-diagonal submatrix is zero. More generally, H will always be block diagonal (or a symmetric permutation thereof) if the base has more than one child. Example (b) shows the opposite extreme: the tree has no branches, and so there is no branch-induced sparsity. Example (c) shows a more typical sparsity pattern; and examples (d) and (e) illustrate the effect of the numbering scheme: two different numberings of the same graph produce matrices that are symmetric permutations of each other.

0 1 2 3 4 5 6 7 H L T L Figure 6.3: Preservation of sparsity in the factorization H = LTL In principle, these two matrices have the same underlying sparsity pattern, and the choice of numbering does not affect our ability to exploit it. It is often the case that branch-induced sparsity can affect a substantial proportion of the elements of H. We can get a rough estimate of its beneficial effect using the following rule of thumb: if H has a density of ρ, then the cost of calculating it is roughly ρ times the cost of calculating a dense H of the same size, and the cost of factorizing it is roughly ρ2 times the cost of factorizing a dense matrix of the same size. ρ is defined to be the proportion of elements in H that are not branch-induced zeros, so 0 < ρ ≤1. It is not unusual to encounter densities of around 0.5, and densities close to zero are possible. Overall, the effect of branch-induced sparsity is to make inertia-matrix methods more efficient on branched kinematic trees than they are on unbranched trees of the same size. The composite-rigid-body algorithm automatically exploits branch-induced sparsity, simply by calculating only the nonzero submatrices of H. However, if we were to try and factorize the resulting matrix using a standard factorization algorithm, then it would treat the matrix as dense, and perform O(n3) arithmetic operations. Therefore, in order to fully exploit the sparsity, we need a factorization algorithm for matrices containing branch-induced sparsity. This turns out to be an easy problem to solve, the solution being to factorize H into either LTL or LTDL, and design the factorization algorithm to skip over the branch-induced zeros. In the sparse matrix literature, the factorization H = LTL would be described as a reordered Cholesky factorization, meaning that it is equivalent to performing a standard Cholesky factorization on a permutation of the original matrix (George and Liu, 1981). Likewise, the factorization H = LTDL would be described as a reordered LDLT factorization. This implies that the LTL and LTDL factorizations have the same numerical properties as the standard Cholesky and LDLT factorizations. The special property of an LTL or LTDL factorization, when applied to a matrix containing branch-induced sparsity, is that the factorization proceeds without fill-in. In other words, every branch-induced zero element in the matrix remains zero throughout the factorization process. (This is proved in Featherstone (2005).) A factorization with this property is said to be optimal (Duff

LTL factorization: for k = n to 1 do Hkk = √ Hkk i = λ(k) while i̸ = 0 do Hki = Hki/Hkk i = λ(i) end i = λ(k) while i̸ = 0 do j = i while j̸ = 0 do Hij = Hij −Hki Hkj j = λ(j) end i = λ(i) end end LTDL factorization: for k = n to 1 do i = λ(k) while i̸ = 0 do a = Hki/Hkk j = i while j̸ = 0 do Hij = Hij −a Hkj j = λ(j) end Hki = a i = λ(i) end end Table 6.3: LTL and LTDL factorization algorithms et al., 1986). One immediate consequence is that the sparsity pattern of H is preserved in the resulting factors, which are therefore also sparse. Figure 6.3 illustrates this effect, showing both the original pattern and the pattern in the resulting triangular factors. This sparsity can be exploited during backsubstitution. A second consequence of the no-fill-in property is that the branchinduced zeros can be ignored completely during the factorization process. 6.5 Sparse Factorization Algorithms Table 6.3 presents the pseudocode for the sparse LTL and LTDL factorization algorithms. These algorithms expect to be given two inputs: an n × n matrix H and an n-element integer array λ. H is expected to be symmetric and positive definite, and the elements of λ are expected to satisfy 0 ≤λ(i) < i. In addition, H is expected to have the following sparsity pattern, which is optimally exploited by the algorithms: On each row k of H, the nonzero elements below the main diagonal appear only in columns λ(k), λ(λ(k)), and so on. This rule identifies the set of elements that the algorithms will treat as being nonzero. All other elements are assumed to be zero, and are ignored. If λ(k) = k −1 for all k, then the algorithms treat every element as nonzero, in which case they perform a dense-matrix LTL or LTDL factorization. Figure 6.4 shows how the algorithms work. They both operate in situ on the given matrix, and never access any element above the main diagonal. The LTL

3 3 3 2 2 1 3 3 3 2 2 1 3 3 3 2 2 1 never accessed finished 2 3 1. 2. 3. LTL LTDL Hkk = Hkk Hki = Hki/Hkk Hij = Hij − Hki Hkj 2a. 2b. 3. H’ki = Hki/Hkk Hki = H’ki Hij = Hij − H’ki Hkj k = 7 k = 6 k = 5 1 k 1. do nothing Figure 6.4: Illustration of the factorization process algorithm computes L and returns it in the lower triangle of H. The LTDL algorithm computes D and L, returns D on the main diagonal, and returns the off-diagonal elements of L below it. This works because the algorithm computes a unit lower-triangular matrix, so its diagonal elements are known to have the value 1, and therefore do not need to be returned. Each algorithm has an outer loop that visits each row in turn, working from n back to 1. At any stage in the factorization process, rows k + 1 to n are finished, and contain rows k + 1 to n of the returned factors. Rows 1 to k can be divided into three areas, as shown in the diagram. Area 1 consists of just the element Hkk; area 2 consists of elements 1 to k −1 of row k; and area 3 consists of the triangular region from rows 1 to k −1. A different calculation takes place in each area, as shown in the figure. The pseudocode for the LTL factorization performs the area-2 and area-3 calculations in two separate loops; but the pseudocode for the LTDL factorization interleaves these calculations in such a way that it never needs to remember more than one H ′ ki value at any point. The current remembered value is stored in the local variable a. The inner loops are designed to iterate only over the values λ(k), λ(λ(k)), and so on. This is where the sparsity is exploited. In effect, the algorithms know where the zeros are, and simply skip over them. Figure 6.4 illustrates the cost reduction by showing the first three steps in the factorization of the matrix H from Figure 6.3. At k = 7, the algorithms perform two lots of the area-2 calculation and three lots of the area-3 calculation; and similarly at k = 6 and k = 5. This is far fewer than the number of calculations that would have been necessary if there had been no zeros.

for i = 1 to n do λ′(i) = i −1 end map(0) = 0 for i = 1 to NB do map(i) = map(i −1) + ni end for i = 1 to NB do λ′(map(i −1) + 1) = map(λ(i)) end 0 1 2 3 4 5 6 0 1 2 3 4 5 6 7 8 Table 6.4: Algorithm for calculating the expanded parent array, and illustration of the expanded connectivity graph Expanding the Parent Array The parent array for a kinematic tree contains NB elements, but the factorization algorithms require an array with n elements. If n = NB then no action is required. Otherwise, an expanded parent array must be constructed for use in the factorization process. Table 6.4 presents an algorithm to do this, and an illustration of what the algorithm does. Essentially, we start with the original connectivity graph, and replace each multi-freedom joint with a chain of single-freedom joints. This produces an expanded graph in which every arc represents a single joint variable. This new graph is renumbered so that arc i represents the ith variable in ¨q; and the expanded parent array is computed from this graph. The illustration shows the original and expanded connectivity graphs for a kinematic tree in which joint 2 has three degrees of freedom, but every other joint has only one. Arc 2 is therefore replaced by a chain of three arcs, and the new tree is renumbered so that each arc number equals the index of its variable in ¨q. As the variables for joint 2 occupy the second, third and fourth places in ¨q, the new arcs must be numbered 2, 3 and 4. The parent array for the original graph is (0, 1, 1, 2, 2, 3); and the parent array for the expanded graph is (0, 1, 2, 3, 1, 4, 4, 5). The algorithm in Table 6.4 calculates an expanded parent array, λ′, directly from the original. The first step is to initialize λ′ to the array (0, 1, 2, . . ., n−1). This is a clever trick that correctly initializes n −NB of the elements in λ′, leaving only NB elements to be computed from λ. The second loop constructs a mapping array such that map(i) = Pi j=1 nj, where nj is the degree of freedom of joint j. (The variables of joint i appear in ¨q at locations map(i −1) + 1 to map(i) inclusive.) The third loop performs two mappings rolled into one: the value stored in λ(i) is converted to map(λ(i)), which is the correct value to be inserted into λ′; and this value is inserted into λ′ at location map(i −1) + 1, which is the correct place to put it. For the example shown in the table, the mapping array is (0, 1, 4, 5, 6, 7, 8) (indexed from zero); the two entries in λ′

y = L x for i = n to 1 do yi = Lii xi j = λ(i) while j̸ = 0 do yi = yi + Lij xj j = λ(j) end end x = L−1x for i = 1 to n do j = λ(i) while j̸ = 0 do xi = xi −Lij xj j = λ(j) end xi = xi/Lii end y = LTx for i = 1 to n do yi = Lii xi j = λ(i) while j̸ = 0 do yj = yj + Lij xi j = λ(j) end end x = L−Tx for i = n to 1 do xi = xi/Lii j = λ(i) while j̸ = 0 do xj = xj −Lij xi j = λ(j) end end y = Hx for i = 1 to n do yi = Hii xi end for i = n to 1 do j = λ(i) while j̸ = 0 do yi = yi + Hij xj yj = yj + Hij xi j = λ(j) end end Table 6.5: Multiplication and back-substitution algorithms that are correctly initialized in the first loop are λ′(3) and λ′(4); the elements λ(1) · · · λ(3) are mapped unaltered to λ′(1), λ′(2) and λ′(5); and λ(4) · · · λ(6) are incremented by 2 and mapped to λ′(6) · · · λ′(8). Having computed the expanded parent array once, it should be stored in the system model for future use by the factorization algorithms. Back-Substitution The sparsity pattern in H appears also in L; so it is possible to exploit the sparsity in L using the same techniques as for H. Table 6.5 presents algorithms to calculate the quantities Lx, LTx, L−1x and L−Tx for a general vector x, given only λ, L and x. Back-substitution makes use of L−1x and L−Tx. The table also presents an algorithm to calculate Hx. In every case, the sparsity is exploited by having the inner loop iterate over only the nonzero elements of L or H. The algorithms assume that L is a general lower-triangular factor, such as that produced by the LTL factorization. To modify them to work with a unit lower-triangular factor, simply replace Lii with 1. The algorithms that calculate L x LTx and Hx place the result in a separate vector, y, and leave x unaltered. The L x algorithm allows y to be the same vector as x, in which case it works in situ; but the other two require y̸ = x. The algorithms that calculate L−1x and L−Tx both work in situ.

LTL LTDL factorization n√+ D1d + D2(m + a) D1d + D2(m + a) back-substitution 2nd + 2D1(m + a) nd + 2D1(m + a) L x, LTx nm + D1(m + a) D1(m + a) L−1x, L−Tx nd + D1(m + a) D1(m + a) Hx nm + 2D1(m + a) Table 6.6: Computational cost of sparse algorithms To implement y = L−1x or y = L−Tx, simply copy x to y and apply the algorithm to y. Computational Cost Table 6.6 presents computational cost figures for each algorithm. The symbols m, a, d and √denote the cost of a floating-point multiplication, addition, division and square-root calculation, respectively. The quantities D1 and D2 depend on the sparsity pattern in H, and are given by D1 = n X k=1 (dk −1) (6.26) and D2 = n X k=1 dk(dk −1) 2 , (6.27) where dk is the number of joints on the path between body k and the base (i.e., dk = |κ(k)|), and can be regarded as the depth of body k in the tree. The expressions dk −1 and dk(dk −1)/2 count the number of times that area-2 and area-3 calculations are performed during the processing of row k; so D1 and D2 count the total number of area-2 and area-3 calculations, respectively. D1 is also the total number of nonzero elements below the main diagonal. Thus, the number of nonzero elements in H is n + 2D1. If there is no branch-induced sparsity, then D1 and D2 take their maximum possible values, which are D1 = (n2 −n)/2 and D2 = (n3 −n)/6 . (6.28) Plugging these values into the formulae in Table 6.6 produces cost figures that are identical to those of the standard Cholesky and LDLT factorization algorithms. If d denotes the depth of the tree, defined as the depth of the deepest body, then D1 and D2 are bounded by D1 ≤n(d −1) and D2 ≤nd(d −1)/2 . (6.29) These bounds show that the factorization process is at worst O(nd2).

\section{dditional Notes One of the earliest examples of a dynamics algorithm is due to Uicker (1965, 1967). It combined the 4 × 4 matrix representation of kinematics developed by Denavit and Hartenberg (1955) with a 4 × 4 pseudoinertia matrix, and used Lagrangian techniques to obtain the equation of motion for a closed-loop mechanism. At about the same time, Hooker and Margulies (1965) used vectorial (Newton-Euler) dynamics and augmented bodies to obtain the equations of motion for a free-floating kinematic tree representing a satellite. By the mid 1970s, there were already several computer programs available to calculate the forward dynamics of various kinds of rigid-body system; and an excellent review of these early works can be found in Paul (1975). One early idea that proved to be particularly successful is the sparse-matrix formulation used in the program ADAMS (Orlandea et al., 1977). The earliest textbook on computational rigid-body dynamics appeared at this time (Wittenburg, 1977). It presented a detailed method for constructing computer-oriented models of general rigidbody systems, and for calculating their dynamics. The composite-rigid-body algorithm first appeared as method 3 in Walker and Orin (1982). It provided a new method for calculating the joint-space inertia matrix of a kinematic tree, which was substantially faster than previous methods. More efficient versions of this algorithm, and some variations on the basic idea, appeared subsequently in Featherstone (1984, 1987); Balafoutis and Patel (1989, 1991); Lilly and Orin (1991); Lilly (1993); McMillan and Orin (1998). However, the phenomenon of branch-induced sparsity was not recognized until Featherstone (2005). The first example of a propagation method for forward dynamics appeared in Vereshchagin (1974). In this paper, the dynamics was formulated using Gauss’ principle of least constraint, and the resulting minimization solved via dynamic programming. The algorithm in this paper is practically the same as the articulated-body algorithm, but its significance was not recognized until after the articulated-body algorithm had been published. The next example of a propagation algorithm appeared in Armstrong (1979), followed by the preliminary version of the articulated-body algorithm in Featherstone (1983), and then the complete version in Featherstone (1984, 1987). More efficient versions then appeared in Brandl et al. (1988); McMillan and Orin (1995); McMillan et al. (1995). A few other algorithms have turned out to be practically the same as the articulated-body algorithm; for example, the forward dynamics algorithm of Rodriguez et al. (Rodriguez, 1987; Rodriguez et al., 1991) and the O(n) algorithm in Rosenthal (1990). Superficially, it would seem that the inertia-matrix methods and propagation methods are very different. However, two interesting connections have been found. Drawing on Kalman filter theory, Rodriguez et al. present two factorizations of the joint-space inertia matrix: one implies the recursive Newton-Euler algorithm, while the other has an inverse that implies the articulated-body algorithm (Rodriguez, 1987; Rodriguez et al., 1991). The other connection was}
\label{sec:6.6}
found by Ascher et al. (1997). Starting with an equation similar in form to Eq. 3.18, they show that the joint-space inertia matrix can be obtained by applying Gaussian elimination to one permutation of the coefficient matrix, while the articulated-body algorithm arises from applying Gaussian elimination to another permutation of this matrix. Strictly speaking, inertia-matrix and propagation methods are only two out of three possible strategies for calculating the forward dynamics of a kinematic tree. The third possibility is to assemble the equations of motion of the individual bodies, together with the equations of constraint, to form a single large matrix equation (as shown in Chapter 3, for example). This strategy is very useful for closed-loop systems, but it is not competitive with inertia-matrix and propagation methods when applied to kinematic trees. An example of this class of algorithm can be found in Baraff(1996).
